\subsection{\texorpdfstring{* Measurable linear mappings on a Banach space \(^{1}\)}{Measurable linear mappings on a Banach space \^{}\{1\}}}

PROPOSITION 2. --- Let E be a Banach space, F a locally convex space and u a linear mapping from E into F. Assume that for every closed subset B of F, every compact subset X of E and every measure \(\mu\) on X, the intersection \(X \cap u^{-1}(B)\) is \(\mu\) -measurable. Then u is continuous.

First assume that F is the base field. For every compact subset X of E and every measure \(\mu\) on X, the restriction of u to X is \(\mu\) -measurable (INT, IV). Suppose that u is not continuous. Then we can find a sequence of points \((x_{n})\) in E such that \(\sum_{n} \|x_{n}\| < \infty\) and \(|u(x_{n})| \geqslant n\) for every integer n.~Consider the mapping \(g:(t_{n}) \mapsto \sum_{n} t_{n} x_{n}\) from the cube \(C = [0, 1]^{N}\) into E; it is clear that g is continuous. Hence \(f = u \circ g\) is measurable for every measure on C (INT, V); in particular for the measure \(\mu\) which is the product of Lebesgue measures on the factors of C. Hence there exists a compact subset D of C such that \(\mu(D) > \frac{1}{2}\) and such that the restriction of f to D is continuous, hence also bounded. Let M be the upper bound of \(|f|\) on D and let \(p \in N\) be such that \(p \geqslant 4M\) . Let \(s = (s_{n})\) and \(t = (t_{n})\) be two points of D such that \(s_{n} = t_{n}\) for all \(n \neq p\) . Then

\[
f (s) - f (t) = u (\sum_ {n} s _ {n} x _ {n} - \sum_ {n} t _ {n} x _ {n}) = (s _ {p} - t _ {p}) u (x _ {p}).
\]

Since \(|f(s) - f(t)| \leqslant 2\mathbf{M}\) and \(|u(x_p)| \geqslant p \leqslant 4\mathbf{M}\) , we get

\[
\left| s _ {p} - t _ {p} \right| \leqslant \frac {1}{2}.
\]

The Lebesgue-Fubini theorem (INT, V, 2nd ed., § 8, No.~3, cor. 2 of prop. 7) implies that \(\mu(D) \leqslant \frac{1}{2}\) ; this gives a contradiction. Hence \(u\) is continuous.

In the general case, for every \(v \in F'\) , the linear form \(v \circ u\) is continuous, by the preceding argument. Let \((x_{n})_{n \in \mathbb{N}}\) be a sequence of points of E tending to 0; then the sequence \((u(x_{n}))_{n \in \mathbb{N}}\) tends to 0 in F, if F is assigned the topology \(\sigma(\mathrm{F}, \mathrm{F}')\) ; hence this sequence is bounded for \(\sigma(\mathrm{F}, \mathrm{F}')\) and so it is bounded in F (III, p.~27, cor. 3). Since E is bornological (III, p.~12, prop. 2); the linear mapping \(u: E \to F\) is continuous.*

1 The results of this section depend on the book of Integration.

